\section{A finite full-incidence calculus}
\label{sec:incidence-calculus}

The retained-pair question is finite at every value of \(X\).  We therefore
begin with a finite-dimensional theorem which keeps all four pre-terminal
labels.  The theorem distinguishes preservation of the pre-terminal raw
form, preservation of the terminal atomic form, and preservation of the
grouped physical operator.  These notions coincide only under additional
incidence hypotheses.

\subsection{The ambient incidence and its pair projection}

Let \(\cM,\cR\subset\N_{\ge1}\) and \(\cJ\subset\Z\) be finite, and let
\(\Gamma\) be a finite label set.  Put
\begin{equation}
 \Omega:=\cM\times\cR\times\cJ\times\Gamma
 \label{eq:incidence-ambient-product}
\end{equation}
and write
\[
 u=(m,r,j,\gamma),\qquad m(u)=m,
\]
and take the row-coefficient space to be \(Z:=\C^{\cM}\).  All sets,
structural coefficients, and maps below are fixed before \(z\in Z\) is
chosen.  A version with a finite base fiber
\(B_m\subseteq\cJ\times\Gamma\) depending on \(m\) is obtained by replacing
\(\cJ\times\Gamma\) by \(B_m\); none of the proofs changes.

Let \(c=(c_u)_{u\in\Omega}\) be the coefficient-free structural field.  It
contains every literal source shape, mask, opposite-row factor, and orbit
factor preceding terminalization, but not the free row coefficient \(z_m\).
For \(V\subseteq\Omega\), define the zero-extended synthesis operator
\begin{equation}
 (C_Vz)_u:=\one_V(u)c_u z_{m(u)},
 \qquad C_V:Z\longrightarrow\ell^2(\Omega).
 \label{eq:incidence-CV}
\end{equation}
Thus every carrier is compared in the same counting-measure Hilbert space.

\begin{definition}[Full incidence, pair projection, and rectangle]
\label{def:incidence-full-pair-rectangle}
Let \(\cU\subseteq\Omega\) be a declared full pre-terminal incidence.  Its
pair projection and rectangular completion are
\begin{align}
 I(\cU)
 &:=\pr_{m,r}\cU
   =\{(m,r):\text{there exists }(j,\gamma)
                    \text{ with }(m,r,j,\gamma)\in\cU\},
 \label{eq:incidence-pair-projection}\\
 \Rect(I)
 &:=\{(m,r,j,\gamma)\in\Omega:(m,r)\in I\}.
 \label{eq:incidence-rectangular-completion}
\end{align}
The associated spurious-cross set is
\begin{equation}
 \cS(\cU):=\Rect(I(\cU))\setminus\cU.
 \label{eq:incidence-spurious-set}
\end{equation}
\end{definition}

By construction, \(\cU\subseteq\Rect(I(\cU))\).  This containment does not
say that the new crossed coordinates in \(\cS(\cU)\) have zero structural
coefficient or fail the terminal test.

\subsection{Terminal restriction and the grouped operator}

Fix \(h\in\Z\).  For \(u=(m,r,j,\gamma)\), put
\begin{equation}
 p_h(u):=\frac{mj+h}{r}
 \label{eq:incidence-prime-candidate}
\end{equation}
when the quotient is a positive integer.  Let
\(\one_{\mathrm{term},h}(u)\) be zero when that quotient is not a positive
integer, and otherwise the indicator of the complete declared terminal
predicate: the quotient in \eqref{eq:incidence-prime-candidate} is a prime
in the stated support, and every required terminal cutoff,
squarefree condition, coprimality condition, literal mask, and nonvanishing
terminal profile holds.  Define
\begin{equation}
 t_h(u):=\mu(r)\Lambda_X(mj+h)\one_{\mathrm{term},h}(u),
 \qquad
 (D_hx)_u:=t_h(u)x_u.
 \label{eq:incidence-terminal-diagonal}
\end{equation}
Here \(\mu(r)\) is the single inherited cofactor sign, and
\(\Lambda_X(n)=(\log n)\Phi_0(n/X)\) for \(n>0\), zero-extended to
\(\Z\), is the inherited smooth terminal weight rather than the von
Mangoldt function.  No additional M\"obius factor may be inserted in the
physical specialization.

Let \(\mathcal Q\) be a finite output-key set, let
\(q:\Omega\to\mathcal Q\) be fixed, and define
\begin{equation}
 (Gx)_a:=\sum_{\substack{u\in\Omega\\q(u)=a}}x_u,
 \qquad
 G:\ell^2(\Omega)\longrightarrow\ell^2(\mathcal Q).
 \label{eq:incidence-grouping-map}
\end{equation}
For the fixed-shift residual packet the literal key is
\begin{equation}
 q(m,r,j,\gamma)=(\gamma,d_mr,j).
 \label{eq:incidence-physical-key}
\end{equation}
For \(V\subseteq\Omega\), write
\begin{align}
 P_V(z)&:=\|C_Vz\|_2^2,
 \label{eq:incidence-preterminal-form}\\
 A_{h,V}(z)&:=\|D_hC_Vz\|_2^2,
 \label{eq:incidence-terminal-atomic-form}\\
 F_{h,V}&:=GD_hC_V.
 \label{eq:incidence-grouped-operator}
\end{align}
The first two forms use raw counting measure.  The last map is coherent on
each output fiber, so its squared norm need not be a sum of atomic squares.

\subsection{Exact rectangularization identities}

\begin{theorem}[Rectangularization difference identities]
\label{thm:incidence-rectangularization-difference}
Let \(\cU\subseteq V\subseteq\Omega\), and put
\(\cS:=V\setminus\cU\).  Then
\begin{align}
 C_V-C_{\cU}&=C_{\cS},
 \label{eq:incidence-C-difference}\\
 D_hC_V-D_hC_{\cU}&=D_hC_{\cS},
 \label{eq:incidence-D-difference}\\
 F_{h,V}-F_{h,\cU}&=GD_hC_{\cS}.
 \label{eq:incidence-F-difference}
\end{align}
Moreover, for every \(z\in Z\),
\begin{align}
 P_V(z)-P_{\cU}(z)
 &=\sum_{u\in\cS}|c_u|^2|z_{m(u)}|^2,
 \label{eq:incidence-P-positive-difference}\\
 A_{h,V}(z)-A_{h,\cU}(z)
 &=\sum_{u\in\cS}|t_h(u)c_u|^2|z_{m(u)}|^2.
 \label{eq:incidence-A-positive-difference}
\end{align}
Finally, the \((a,m)\)-matrix entry of the grouped difference is
\begin{equation}
 \boxed{
 (F_{h,V}-F_{h,\cU})_{a,m}
 =\sum_{\substack{u\in\cS\\q(u)=a,\ m(u)=m}}t_h(u)c_u.}
 \label{eq:incidence-grouped-matrix-entry}
\end{equation}
In particular, the identities apply to
\(V=\Rect(I(\cU))\) and \(\cS=\cS(\cU)\).
\end{theorem}

\begin{proof}
Since \(V=\cU\,\dot\cup\,\cS\), one has
\(\one_V-\one_{\cU}=\one_{\cS}\).  Substitution in
\eqref{eq:incidence-CV} proves \eqref{eq:incidence-C-difference};
applying \(D_h\) and \(G\) gives the next two identities.  The vectors on
\(\cU\) and \(\cS\) have disjoint coordinate supports, before and after
the diagonal map \(D_h\), so Pythagoras proves the two form identities.
Finally, evaluate \eqref{eq:incidence-F-difference} at the row basis vector
\(e_m\) and read its \(a\)-coordinate.
\end{proof}

The positive identities
\eqref{eq:incidence-P-positive-difference} and
\eqref{eq:incidence-A-positive-difference} must not be transferred through
\(G\): terms in \eqref{eq:incidence-grouped-matrix-entry} can cancel.

\begin{theorem}[Three exact fidelity criteria]
\label{thm:incidence-three-fidelity-criteria}
Let \(I=I(\cU)\), and put
\(\cS=\Rect(I)\setminus\cU\).  Then:
\begin{enumerate}[label=\textup{(\roman*)},leftmargin=3em]
 \item The rectangular lift preserves the pre-terminal synthesis and raw
       form exactly if and only if
       \begin{equation}
        \begin{aligned}
        C_{\Rect(I)}=C_{\cU}
        &\ \Longleftrightarrow\
        P_{\Rect(I)}=P_{\cU}\text{ on }Z\\
        &\ \Longleftrightarrow\
        c_u=0\quad(u\in\cS).
        \end{aligned}
        \label{eq:incidence-preterminal-iff}
       \end{equation}
 \item It preserves the terminal atomic synthesis and raw form exactly if
       and only if
       \begin{equation}
        \begin{aligned}
        D_hC_{\Rect(I)}=D_hC_{\cU}
        &\ \Longleftrightarrow\
        A_{h,\Rect(I)}=A_{h,\cU}\text{ on }Z\\
        &\ \Longleftrightarrow\
        t_h(u)c_u=0\quad(u\in\cS).
        \end{aligned}
        \label{eq:incidence-terminal-iff}
       \end{equation}
 \item It preserves the grouped physical operator exactly if and only if
       \begin{equation}
        \boxed{
        \begin{aligned}
        F_{h,\Rect(I)}=F_{h,\cU}
        \quad\Longleftrightarrow\quad&
        \sum_{\substack{u\in\cS\\q(u)=a,\ m(u)=m}}t_h(u)c_u=0\\[-2pt]
        &\text{for every }(a,m).
        \end{aligned}}
        \label{eq:incidence-grouped-iff}
       \end{equation}
\end{enumerate}
Consequently, pre-terminal fidelity implies terminal-atomic fidelity,
which implies grouped-operator fidelity.  Neither converse holds for a
general grouping map.
\end{theorem}

\begin{proof}
Map equality implies equality of the corresponding squared norms.
Conversely, take \(z=e_m\) in
\eqref{eq:incidence-P-positive-difference}.  Equality of the pre-terminal
forms gives
\[
 \sum_{\substack{u\in\cS\\m(u)=m}}|c_u|^2=0
 \qquad(m\in\cM),
\]
so every \(c_u\) on \(\cS\) vanishes.  The terminal argument is identical
using \eqref{eq:incidence-A-positive-difference}.  The grouped criterion
is \eqref{eq:incidence-grouped-matrix-entry}, because a linear map from
\(\C^{\cM}\) is zero exactly when every matrix entry is zero.
\end{proof}

\begin{example}[The hierarchy is strict]
\label{ex:incidence-strict-hierarchy}
If a spurious coordinate \(u_0\) has \(c_{u_0}\ne0\) but
\(t_h(u_0)=0\), pre-terminal fidelity fails while both later fidelities can
hold.  Separately, two spurious atoms in the same row and output fiber with
terminal coefficients \(1\) and \(-1\) violate terminal-atomic fidelity
but cancel in the grouped operator.  The latter is an abstract example;
the literal key \eqref{eq:incidence-physical-key} has the singleton
property below.
\end{example}

\subsection{Spurious-cross witnesses}

\begin{proposition}[Exact spurious-cross witnesses]
\label{prop:incidence-spurious-cross-witness}
Let \(I=I(\cU)\), let \(\cS=\Rect(I)\setminus\cU\), and take
\(u_0\in\cS\), with \(m_0=m(u_0)\) and \(a_0=q(u_0)\).
\begin{enumerate}[label=\textup{(\alph*)},leftmargin=2.7em]
 \item If \(c_{u_0}\ne0\), then
       \begin{equation}
        P_{\Rect(I)}(e_{m_0})-P_{\cU}(e_{m_0})
        \ge |c_{u_0}|^2>0.
        \label{eq:incidence-preterminal-witness}
       \end{equation}
 \item If \(t_h(u_0)c_{u_0}\ne0\), then
       \begin{equation}
        A_{h,\Rect(I)}(e_{m_0})-A_{h,\cU}(e_{m_0})
        \ge |t_h(u_0)c_{u_0}|^2>0.
        \label{eq:incidence-terminal-witness}
       \end{equation}
 \item If
       \begin{equation}
        \Sigma_{a_0,m_0}(\cS)
        :=\sum_{\substack{u\in\cS\\q(u)=a_0,\ m(u)=m_0}}
                    t_h(u)c_u\ne0,
        \label{eq:incidence-spurious-fiber-sum}
       \end{equation}
       then
       \begin{equation}
        \|(F_{h,\Rect(I)}-F_{h,\cU})e_{m_0}\|_2
        \ge |\Sigma_{a_0,m_0}(\cS)|>0.
        \label{eq:incidence-grouped-witness}
       \end{equation}
       In particular, this holds when \(u_0\) is the only terminal-active
       spurious atom in its \((a_0,m_0)\)-fiber.
\end{enumerate}
\end{proposition}

\begin{proof}
The first two assertions are the single-summand lower bounds in
\eqref{eq:incidence-P-positive-difference} and
\eqref{eq:incidence-A-positive-difference}.  For the last assertion, the
\(a_0\)-coordinate of the displayed difference vector is exactly
\(\Sigma_{a_0,m_0}(\cS)\).
\end{proof}

\begin{corollary}[The literal residual key removes accidental rowwise
cancellation]
\label{cor:incidence-physical-key-singleton}
Assume every row integer \(m\) determines \(d_m\), and use
\(q(u)=(\gamma,d_mr,j)\).  For fixed \((a,m)\), at most one
\(u\in\Omega\) satisfies \(q(u)=a\) and \(m(u)=m\).  Hence
\begin{equation}
 F_{h,\Rect(I)}=F_{h,\cU}
 \quad\Longleftrightarrow\quad
 D_hC_{\Rect(I)}=D_hC_{\cU}.
 \label{eq:incidence-physical-grouped-atomic-equivalence}
\end{equation}
Cancellation between different rows cannot repair this uniform equality,
because the row variables \(z_m\) are independent.
\end{corollary}

\begin{proof}
Fix \(a=(\gamma,s,j)\) and \(m\).  The row determines \(d_m\), and
\(s=d_mr\) determines \(r\).  Thus \(m,r,j,\gamma\), hence \(u\), are
fixed.  Every sum in \eqref{eq:incidence-grouped-iff} has at most one term.
\end{proof}

This corollary does not identify the inherited relation \(\cI_X\) with the
projection of any particular full incidence.  It says only that, after the
literal full incidence has been written down, an active spurious cross
cannot be hidden by the residual-key grouping
\(q=(\gamma,d_mr,j)\) in a coefficient-uniform comparison.  A later map
which erases \(d_mr\) can merge distinct cofactors and requires the separate
fiber/kernel test of \cref{sec:tpc63-labeled-completion}.
